\documentclass[11pt,a4paper]{article}
\usepackage[utf8]{inputenc}
\usepackage[T1]{fontenc}
\usepackage{amsfonts}
\usepackage[french]{babel}
\frenchbsetup{StandardLists=true} % à inclure si on utilise \usepackage[francais]{babel}
\usepackage{enumitem}
\usepackage{amssymb}
\usepackage{amsmath}
\usepackage[left=2cm,right=2cm,top=2cm,bottom=2cm]{geometry}
\usepackage{multirow}
\usepackage[dvipsnames]{xcolor}

\usepackage[T1]{fontenc}
\usepackage{fouriernc}
\usepackage{graphicx}

\usepackage{setspace}
\singlespacing

\usepackage{pstricks}
\usepackage{xkeyval}
\usepackage{pst-labo}


\usepackage{multicol}
\usepackage{caption}
\usepackage{fancyhdr}
\pagestyle{fancy}
\renewcommand\headrulewidth{1pt}
\fancyhead[L]{Lycée Denis Diderot }
\fancyhead[R]{Classe de 1 Bac Pro }
\setcounter{page}{1}\fancyfoot[C]{page \thepage}
\fancyfoot[R]{Année 2024-2025}
\fancyfoot[L]{Alexis RUIZ}
\usepackage{multirow,tabularx,booktabs}
\usepackage[tikz]{bclogo} 

\newcommand{\cadreblanc}[1]{%
\medskip\framebox[\linewidth][c]{%
\rule{0mm}{#1mm}}\par
\medskip}

\usepackage{awesomebox}
\setlength{\aweboxvskip}{0mm}
\usepackage{pifont}
\usepackage{chemmacros}
\chemsetup{modules=all}
\setlength{\columnseprule}{.25mm}

\renewcommand{\arraystretch}{2.5}
\newcommand*\acocher{\quitvmode{\fboxsep0pt \fboxrule0.8pt \fbox{\vrule height1.8ex depth.1ex width0pt \vrule height0pt depth0pt width1.9ex }}}

\newcolumntype{C}[1]{>{\centering\arraybackslash }b{#1}}

\newcommand{\code}{\awesomebox[black]{2pt}{

\includegraphics[scale=0.15]{Chapitre 1 - Energie et puissance électrique/images/frame.png}
}{black}}

\newcommand{\moteur}{\awesomebox[black]{2pt}{

\includegraphics[scale=1]{images/moteur.jpg} 
}{black}}

\newcommand{\defconclusion}{\awesomebox[black]{2pt}{\rotatebox{90}{\Large{\textbf{Conclusion}}}}{black}}

\newcommand{\defbox}{\awesomebox[black]{2pt}{
\includegraphics[scale=0.5]{Chapitre 1 - Energie et puissance électrique/images/appel exam.png}}{black}}

  \newenvironment{itemize*}%  % listes compactes
    {\begin{itemize}%
      \setlength{\itemsep}{0pt}%
      \setlength{\parskip}{0pt}}%
    {\end{itemize}}

\frenchbsetup{StandardLists=true}



\begin{document}

\flushleft \textbf{NOM : \hfill Classe : \hspace{3cm} ~~\\[2mm]
Prénom:\\[0,4cm]
}
\hrule\
\\[0,2cm]

\begin{center}
 \LARGE{TP2  - Mesure d'un déphasage  }
\end{center}

\vspace{3 mm}

\normalsize

\hrule\
\\[0,2cm]
   \begin{center}
   \textbf{Le soin et la rédaction interviendront dans la notation.}
   \end{center} 

\fcolorbox{black}{blue!30}{\textbf{Introduction}}
\begin{multicols}{2}
\includegraphics[scale=0.32]{Chapitre 3 - Puissance consommée par un appareil éléctrique/images/bobine-allumage-basse-tension-interieur-ciao-px.jpg}

Dans l’industrie, un grand nombre de systèmes sont
constitués de circuits inductifs. \\ On les retrouve dans les
moteurs, les générateurs, et les transformateurs. Ils
sont également utilisés dans des circuits où ils
réalisent des fonctions plus compliquées.

\columnbreak


\includegraphics[scale=0.32]{Chapitre 3 - Puissance consommée par un appareil éléctrique/images/TP2 -Condo.jpg}

\vspace{7mm}

La plupart des circuits électroniques comportent des
condensateurs. Leurs rôles vont du stockage de
l’énergie au filtrage.\\ Ils sont également utilisés dans
des circuits électroniques où ils réalisent des fonctions
plus compliquées (déphasage ; résonance ).



\end{multicols}

\vfill 
\dotfill 
\vfill

\noindent%
\begin{minipage}{.42\textwidth}%

    
 \begin{bclogo}[couleur=blue!30 ,  arrondi =0.1 , logo = \bcinfo, epBarre =2]{ Matériel }

\begin{itemize*}
    \item 1 GBF
    \item 1 bobine 100 mH (0,1 H). 
    \item 2 multimètres
    \item 1 plaque de test 
    \item 1 résistor de 330 $\mathrm{\Omega}$
    \item 1 condensateur 4,7 $\mathrm{\mu F}$
\end{itemize*}
\end{bclogo}

\end{minipage}%
\hfill
\begin{minipage}{.52\textwidth}%
\includegraphics[width=\textwidth]{Chapitre 3 - Puissance consommée par un appareil éléctrique/images/TP2-montage.jpg}
\end{minipage}%

\vfill 

\fcolorbox{black}{blue!30}{\textbf{1. Visualisation du déphasage entre u et i dans une bobine}}

\vfill 

\fcolorbox{black}{blue!30}{\textbf{Montage}}
\begin{itemize*}
    \item Réaliser le montage schématisé ci-dessus \textbf{(L = 0,1 H)}
    \item Réglages préliminaires : 
    \begin{itemize*}
        \item Régler $\mathrm{U=6V}$ et $f\mathrm{=1 000 Hz}$
        \item Régler l'oscilloscope afin de visualiser une période sur l'écran avec l'amplitude la plus grande. 
    \end{itemize*}
    \end{itemize*}
    \defbox{\fbox{
\begin{minipage}{10cm}{

Appel $\mathrm{n^\circ 1}$ - Appeler M Ruiz pour qu'il vérifie le montage.}

\end{minipage}
}}

\vfill  
\newpage


\fcolorbox{black}{blue!30}{\textbf{Visualisation}}\\
 

\noindent%
\begin{minipage}{.38\textwidth}%
   
 \begin{bclogo}[couleur=blue!30 ,  arrondi =0.1 , logo = \bcinfo, epBarre =2]{ Relever l'oscillogramme }

\begin{itemize*}
    \item \textbf{Voie 1 : } On visualise la tension u aux bornes de la bobine. 
    \vspace{3mm}
    \item \textbf{Voie 2 : }On visualise la tension aux bornes de R (on a ainsi une image du
courant i qui traverse la bobine).
\end{itemize*}
\end{bclogo}


\end{minipage}%
\hfill
\begin{minipage}{.55\linewidth}%
\includegraphics[width=\textwidth]{Chapitre 3 - Puissance consommée par un appareil éléctrique/images/TP2-Oscillo1.jpg }
\end{minipage}%

\vfill 

\fcolorbox{black}{blue!30}{\textbf{Calcul du déphasage \Large{{$\mathrm{\varphi}$}} 
\normalsize 
entre u (voie 1) et i (voie 2) :}}\\[5mm]

\begin{itemize*}
    \item Relever le décalage $x$ (en s) entre les deux courbes puis déterminer \Large{$\mathrm{\varphi}$} à l'aide de la formule : \\[5mm]

\begin{center}
    \LARGE{$\mathrm{\varphi} = \dfrac{x \times 360 }{T}=$  }
\end{center}
\normalsize
   \item Préciser quelle tension est en avance sur l’autre : \\[5mm]

   \dotfill 
   
\end{itemize*}

\vfill 

\fcolorbox{black}{blue!30}{\textbf{2. Visualisation du déphasage entre u et i dans un condensateur}} \\[5mm]

\begin{minipage}{.38\textwidth}%
\fcolorbox{black}{blue!30}{\textbf{Montage}}
\begin{itemize*}
    \item Réaliser le montage schématisé ci-contre \textbf{(C = 4,7 $\mathrm{\mu F}$)}
    \item Réglages préliminaires : 
    \begin{itemize*}
        \item Régler $\mathrm{U=6V}$ et $f\mathrm{=25Hz}$
        \item Régler l'oscilloscope afin de visualiser une période sur l'écran avec l'amplitude la plus grande. 
    \end{itemize*}
    \end{itemize*}
    \end{minipage}
    \hfill  
    \begin{minipage}{.48\textwidth}%
        \includegraphics[width=\textwidth]{Chapitre 3 - Puissance consommée par un appareil éléctrique/images/TP2-montage2.jpg}
    \end{minipage}

    \vfill 
    \defbox{\fbox{
\begin{minipage}{10cm}{

Appel $\mathrm{n^\circ 2}$ - Appeler M Ruiz pour qu'il vérifie le montage.}

\end{minipage}
}}

\vfill  
\newpage


\fcolorbox{black}{blue!30}{\textbf{Visualisation}}\\
 

\noindent%
\begin{minipage}{.38\textwidth}%
   
 \begin{bclogo}[couleur=blue!30 ,  arrondi =0.1 , logo = \bcinfo, epBarre =2]{ Relever l'oscillogramme }

\begin{itemize*}
    \item \textbf{Voie 1 : } On visualise la tension u aux bornes du condensateur. 
    \vspace{3mm}
    \item \textbf{Voie 2 : }On visualise la tension aux bornes de R (on a ainsi une image du
courant i qui traverse le condensateur).
\end{itemize*}
\end{bclogo}


\end{minipage}%
\hfill
\begin{minipage}{.55\linewidth}%
\includegraphics[width=\textwidth]{Chapitre 3 - Puissance consommée par un appareil éléctrique/images/TP2-Oscillo2.jpg }
\end{minipage}%

\vfill 

\fcolorbox{black}{blue!30}{\textbf{Calcul du déphasage \Large{{$\mathrm{\varphi}$}} 
\normalsize 
entre u (voie 1) et i (voie 2) :}}\\[5mm]

\begin{itemize*}
    \item Relever le décalage $x$ (en s) entre les deux courbes puis déterminer \Large{$\mathrm{\varphi}$} \normalsize à l'aide de la formule : \\[5mm]

\begin{center}
    \LARGE{$\mathrm{\varphi} = \dfrac{x \times 360 }{T}=$  }
\end{center}
\normalsize
   \item Préciser quelle tension est en avance sur l’autre : \\[8mm]

   \dotfill 
   
\end{itemize*}


\vfill

\fcolorbox{black}{blue!30}{\textbf{3. Conclusion}} \\[5mm]






\begin{center}
    

    \begin{tabular}{|C{2.2cm}|C{2cm}|C{8.3cm}|C{2cm}|C{2cm}|}
         \hline
          & Symbole & Représentation de Fresnel temporelle & Déphasage \\
          \hline
          Bobine parfaite & \includegraphics[scale=0.3]{Chapitre 3 - Puissance consommée par un appareil éléctrique/images/TP2-bobine1.jpg}  & \includegraphics[scale=0.3]{Chapitre 3 - Puissance consommée par un appareil éléctrique/images/TP2-bobine2.jpg}  & \LARGE{$\mathrm{\varphi} = .... $} \\
          \hline
          Condensateur parfait & \includegraphics[scale=0.3]{Chapitre 3 - Puissance consommée par un appareil éléctrique/images/TP2-condo1.jpg}   & \includegraphics[scale=0.3]{Chapitre 3 - Puissance consommée par un appareil éléctrique/images/TP2-condo2.jpg} & \LARGE{$\mathrm{\varphi} = .... $} \\
          \hline
          
          
        
 
    \end{tabular}
\end{center}

\vfill 

\newpage 





\fcolorbox{black}{blue!30}{\textbf{4. Application - circuit RLC série (étude de la résonance}} \\[5mm]


\begin{minipage}{.38\textwidth}%
\fcolorbox{black}{blue!30}{\textbf{Montage}}
\begin{itemize*}
    \item Réaliser le montage schématisé ci-contre \textbf{(C = 4,7 $\mathrm{\mu F}$ ; L = 100 H) }
    \item Réglages préliminaires : 
    \begin{itemize*}
        \item Régler $\mathrm{U=6V}$ et $f\mathrm{=1000 Hz}$
        \item Régler l'oscilloscope afin de visualiser une période sur l'écran avec l'amplitude la plus grande. 
    \end{itemize*}
    \end{itemize*}
    \end{minipage}
    \hfill  
    \begin{minipage}{.48\textwidth}%
        \includegraphics[width=\textwidth]{Chapitre 3 - Puissance consommée par un appareil éléctrique/images/TP2-montage3.jpg}
    \end{minipage}

    \vfill 
    \defbox{\fbox{
\begin{minipage}{10cm}{

Appel $\mathrm{n^\circ 3}$ - Appeler M Ruiz pour qu'il vérifie le montage.}

\end{minipage}
}}

\vfill


\fcolorbox{black}{blue!30}{\textbf{Mesures}} \\[5mm]
On fait varier la fréquence $f$ et on relève le courant  I qui traverse le circuit. 
\renewcommand{\arraystretch}{3.5}
\begin{tabular}{|C{2.2cm}|C{1.5cm}|C{1.5cm}|C{1.5cm}|C{3cm}|C{1.5cm}|C{1.5cm}|C{1.5cm}|}
\hline
$f$~~(Hz) & 100 & 150 & 200 & $f_R = ~~~~~~$ & 250 & 500 & 1000 \\
\hline
U1 (V) & 6 & 6 & 6 & 6 & 6 & 6 & 6  \\
\hline
U2 (V)  & & & & $\mathrm{U_{2Max}= ~~~~}$ & & & \\
\hline
\end{tabular}


\vfill 

\begin{itemize*}
    \item Que peut-on dire alors du déphasage entre u et i lorsque $f = f_R$? \\[5mm]
    \dotfill  
    \vspace{5mm}
    \item Retrouver la fréquence de la résonance $f_R$ en appliquant la relation : \\[3mm]
    
    \Large{$$f_R=\dfrac{1}{2\pi \sqrt{L \times C}}= ~~~~~~~~~~$$}
\end{itemize*}

\vfill 

\end{document}